% chapters/09/rangkuman.tex
\bagian{Rangkuman}

\begin{rangkumanbox}
\begin{itemize}
  \item Variabel bernilai benar-salah pada Bab~VIII sudah merupakan mesin
        keadaan dua keadaan. Yang ditambahkan mesin keadaan yang sesungguhnya
        adalah keadaan yang lebih dari dua, syarat perpindahan yang tersendiri,
        dan pekerjaan penyiapan ketika sebuah keadaan dimasuki.
  \item NPC pada bab ini memiliki empat keadaan: berdiam, berpatroli, mengejar,
        dan menyerang, dengan tujuh perpindahan yang sah.
  \item Diagram keadaan--transisi memakai tiga lambang, yaitu kotak untuk
        keadaan, anak panah untuk perpindahan, dan keterangan pada anak panah
        untuk syaratnya.
  \item Setiap perpindahan harus memiliki syarat yang dapat diperiksa, setiap
        keadaan harus memiliki jalan keluar, dan tidak boleh ada dua syarat
        yang saling bertentangan.
  \item Tiga jenis pewaktu dipakai: pewaktu berdiam, pewaktu tenggang, dan
        pewaktu jeda serangan. Ketiganya dijalankan dengan cara yang sama, yaitu
        ditambah selang waktu lalu dibandingkan dengan nilai batasnya.
  \item Di dalam kode, tiga pekerjaan dipisahkan: \code{decide} menentukan
        keadaan berikutnya, \code{enter_state} menyiapkan keadaan baru, dan
        \code{act} menjalankan perilakunya. Fungsi \code{decide} tidak mengubah
        apa pun, sehingga keputusan dapat diperiksa terpisah dari tindakan.
  \item Mesin keadaan diuji dengan tiga cara: mencatat perpindahan beserta
        bingkainya, menjalankan skenario yang dapat diulang, dan memeriksa
        keabsahan setiap perpindahan terhadap tabel perpindahan yang sah.
  \item Pemeriksaan keabsahan merupakan pemeriksaan sifat yang harus selalu
        berlaku; ia tidak menuntut urutan tertentu, sehingga tahan terhadap
        perbedaan kecepatan komputer.
  \item Ketujuh perpindahan pada skenario pengujian terjadi seluruhnya, dan
        jumlah serangan yang tercatat, yaitu lima, sesuai dengan hitungan dari
        lama keadaan menyerang dan pewaktu jedanya.
  \item Apabila hasil pengujian tidak sesuai dengan dugaan, periksalah lebih
        dahulu apakah dugaan atau rancangan pengujiannya yang keliru.
  \item Mesin keadaan menentukan niat, bukan rute. NPC pada bab ini masih belum
        dapat mengelilingi rintangan; kemampuan itu ditambahkan pada Bab~X
        melalui pencarian jalur.
\end{itemize}
\end{rangkumanbox}
